This chapter examines the security risks addressed by Fiorino, the protections
implemented for them, and the limits that remain.

\paragraph{Setting and intrinsic limitations}
The intended setting is a
local user protecting files from other users of the same machine.

Fiorino essentially protects files behind two factors: access to the encrypted
files (i.e. compromise of the OS, the filesystem, or the system user) and password
for the encryption of the files themselves.

\section{Path handling and filesystem attacks}

Vault and output names are externally supplied and thus untrusted; without
protections, \textit{path traversal} attacks and not only would be possible.

\paragraph{Vault and filename rejection}
Fiorino treats names as single path
components and rejects \texttt{/}, \texttt{.}, and \texttt{..}. On POSIX,
backslash is an ordinary filename character: the implementation accepts it,
but if ever porting this software to Windows the developer must be aware of it.
Constructed paths are length-checked before allocation or use.

\paragraph{Vault directories and \texttt{.auth}}
Vault directories are restricted to the owner, and the code checks their ownership
and permissions.

\paragraph{Symlink rejection}
Fiorino checks vault
directories with \texttt{lstat}, opens sensitive files with
\texttt{O\_NOFOLLOW}, and inspects opened descriptors with \texttt{fstat}.
Listing uses \texttt{fstatat} with \texttt{AT\_SYMLINK\_NOFOLLOW} and such
options are used all across the library.

\section{Password and encrypted-file attacks}

The password verifier and each encrypted file use independent random salts,
while the password is the same across each vault.

\paragraph{Key derivation}
PBKDF2-HMAC-SHA-256 is configured for 600,000 iterations, and verifier
comparison uses \texttt{CRYPTO\_memcmp}.

It does not prevent offline guessing if an attacker obtains a ciphertext,
thus password strength is still relevant (a rate limit is not implemented in the CLI
as it would be useless) while reversal stays computationally unfeasible.

\paragraph{Encryption and authenticity at rest}
AES-256-GCM detects modification of the encrypted payload. Decryption checks
the tag before opening the plaintext destination, so a wrong password or
modified bundle does not truncate an existing output file.

\paragraph{Limitations and future possibilities for AEAD}
The bundle does
not authenticate the vault name or filename as associated data. Consequently,
those names are not cryptographically bound to the ciphertext; a bundle can
be renamed or copied, and the filename and approximate plaintext size remain
visible. The current format also has no version or magic field. Future versions
can use AEAD (Authenticated Encryption with Associated Data) easily on top.

\section{Malformed input and resource limits}

\paragraph{Malformed input}
Malformed or very large inputs can cause out-of-bounds access or excessive
memory use if lengths are trusted. Bundle deserialization checks that the
header and tag are present, checks the ciphertext against the 10,000,000-byte
limit, and allocates only after those checks.

The serialized format is parsed field by field rather than by casting file
bytes to a C structure.

\paragraph{Limitation of input size}
Being AES a block cipher, it introduces padding: this means that there should
be a tolerance in the size of the ciphertext on top of that of the maximum accepted
plaintext; currently, this is not implemented.

\paragraph{Limitation of input from stream}
Taking input from a stream presents some risks that are not fully handled by
the internals of the API, mainly because of the assumption of such stream not
coming from a network socket or equivalently from a slowly yielding pipe.

\section{File consistency and failure behavior}

Encryption takes a shared \texttt{flock}, checks the open stream with
\texttt{fstat}, and reads progressively. It detects a short read and checks
for an extra byte after the measured size. This catches some changes during
the read.

\paragraph{Limitations of \texttt{flock}}
By design \texttt{flock} is advisory: a writer that ignores the lock can
still modify the file, including without changing its size, thus eluding the checks.

\paragraph{Writing ciphertext}
Encrypted output is written to a temporary file, synchronized, renamed (atomically)
over the destination, and followed by a directory synchronization.

\paragraph{Known limitation on writing the ciphertext}
A write failure before the rename can leave the
temporary file behind. Decryption writes
plaintext directly to its destination; an I/O failure during that write can
leave a partial plaintext file. These are known limits of the current file
handling.

\section{Deletion and memory exposure}

Deletion uses \texttt{unlink}, which removes a directory entry; it does not
securely erase data from the storage device. Filesystems, snapshots, copy-on-write
storage, and solid-state devices make overwriting a file an unreliable general
sanitization method. If \texttt{stat} succeeds and reports more than one hard
link, Fiorino refuses deletion and tells the CLI user why.

\paragraph{Known limits and future work}
The link-count check can still race with a concurrent process
that changes the links after the check. Fiorino does not yet enumerate the other
link names.

\paragraph{Memory cleansing in process memory}
The program cleanses sensitive buffers with \texttt{OPENSSL\_cleanse} during
cleanup, and clears unauthenticated plaintext after a GCM verification
failure. This reduces the lifetime of those particular buffers. It cannot
remove copies held in registers, OpenSSL internals, swap, crash dumps, or
other process memory.

\section{Security testing}

The tests exercise concrete rejection and integrity properties, rather than
trying to establish that Fiorino is free of vulnerabilities:

Manual testing previously showed that using a symlink as a vault directory
could redirect encrypted output outside the Fiorino root. The vault checks
were tightened, and the regression test now requires the operation to fail
without creating the authentication record or ciphertext in the symlink
target.

\begin{itemize}
    \item \texttt{tests/fiorino\_roundtrip.c} checks binary and empty-file
    round trips, wrong passwords, truncated bundles, and changes to the IV,
    salt, ciphertext, and tag. It also checks that failed authentication
    preserves an existing plaintext file.
    \item \texttt{tests/fiorino\_limits.c} checks overlong vault and output
    paths and an input larger than the configured limit.
    \item \texttt{tests/fiorino\_vault\_root.c} and
    \texttt{tests/fiorino\_vault\_symlink.c} check rejection of symlinked
    vault roots and vault directories. The latter confirms that a symlinked
    vault does not create the authentication record or ciphertext in its
    target.
    \item \texttt{tests/fiorino\_symlink.c} checks that symlinked plaintext
    outputs and temporary encrypted outputs do not redirect writes.
    \item \texttt{tests/fiorino\_auth\_record.c} checks authentication
    records that are symlinks, FIFOs, or owned by another user when the test
    environment permits the ownership change.
    \item \texttt{tests/encrypt\_file\_cleanup.c} checks encryption success,
    write failure, and rename failure. It manually removes the temporary file
    left by the write-failure case; the test does not claim that the library
    removes it automatically.
    \item CLI tests check usage errors, encryption/decryption round trips,
    listing, deletion, wrong passwords, hard links, and symlink cases. The
    rejection script also supplies hostile-looking vault values, but invokes
    the commands without the required filename, so those values do not reach
    the library's path validation.
\end{itemize}

\paragraph{Functional, sanitization and memory testing}
The default build enables AddressSanitizer and UndefinedBehaviorSanitizer.
\texttt{make test} runs the functional and rejection tests, while
\texttt{make memtest} rebuilds the tests without sanitizers and runs them
under Valgrind.
